% =====================================================================
%  FICHE 8 — THÈME 5 — NIVEAU DIFFICILE — CORRIGÉ
% =====================================================================
\documentclass[11pt,a4paper]{article}
\usepackage[utf8]{inputenc}
\usepackage[T1]{fontenc}
\usepackage{lmodern}
\usepackage[french]{babel}
\usepackage{amsmath,amssymb,mathtools}
\usepackage{icomma}
\usepackage{siunitx}
\sisetup{output-decimal-marker={,},per-mode=power,group-separator={\,},group-minimum-digits=5}
\usepackage[version=4]{mhchem}
\usepackage{xcolor}
\usepackage{colortbl}
\usepackage{tikz}
\usepackage[most]{tcolorbox}
\usepackage{enumitem}
\usepackage{array}
\usepackage{booktabs}
\usepackage{fancyhdr}
\usepackage[hidelinks]{hyperref}
\usetikzlibrary{arrows.meta,positioning,calc}
\usepackage[a4paper,margin=2.1cm,headheight=26pt,headsep=10pt]{geometry}
\usepackage{titlesec}

\definecolor{primary}{RGB}{150,98,162}
\definecolor{primarydark}{RGB}{92,52,110}
\definecolor{excol}{RGB}{0,135,90}
\definecolor{attncol}{RGB}{200,40,55}

\tcbset{boxbase/.style={enhanced,breakable,boxrule=0.9pt,arc=2.5pt,left=3mm,right=3mm,top=2.3mm,bottom=2.3mm,fonttitle=\bfseries,coltitle=white,toptitle=0.6mm,bottomtitle=0.6mm}}
\newtcolorbox[auto counter]{correction}[1][]{boxbase,colback=excol!4,colframe=excol,title={\textsc{Corrigé}~\thetcbcounter\ \ifx\relax#1\relax\relax\else\textmd{--- \itshape #1}\fi}}

\newcommand{\rep}[1]{\textcolor{primarydark}{\textbf{#1}}}
\newcommand{\oui}{\textcolor{excol}{\textbf{Oui}}}
\newcommand{\non}{\textcolor{attncol}{\textbf{Non}}}

\pagestyle{fancy}
\fancyhf{}
\fancyhead[L]{\small\textcolor{primarydark}{\textbf{Fiche 8} — Thème 5 : Équilibre chimique et Le Chatelier}}
\fancyhead[R]{\small\textcolor{primarydark}{\textsc{Difficile — corrigé}}}
\fancyfoot[C]{\small\thepage}
\renewcommand{\headrulewidth}{0.8pt}
\renewcommand{\headrule}{\hbox to\headwidth{\color{primary}\leaders\hrule height \headrulewidth\hfill}}
\setlength{\parindent}{0pt}\setlength{\parskip}{3pt}

\begin{document}

\begin{center}
\begin{tikzpicture}
\node[rectangle,rounded corners=7pt,fill=excol,text=white,minimum width=16cm,
      minimum height=1.1cm,align=center,inner sep=6pt]
  {{\large\bfseries Thème 5 — Équilibre chimique et Le Chatelier}\\[1pt]{\small Niveau \textsc{Difficile} — corrigé détaillé}};
\end{tikzpicture}
\end{center}

\begin{correction}[tableau d'avancement et équation du second degré]
\begin{enumerate}[label=\alph*),leftmargin=2em,topsep=2pt,itemsep=2pt]
  \item Tableau des concentrations (\si{\mol\per\litre}) :
  \begin{center}\small
  \begin{tabular}{@{}lccc@{}}
  \toprule
   & $\ce{N2O4}$ & $\rightleftharpoons$ & $2\,\ce{NO2}$ \\
  \midrule
  initial   & $1{,}0$ & & $0$ \\
  équilibre & $1-x$   & & $2x$ \\
  \bottomrule
  \end{tabular}
  \end{center}
  \item $K = \dfrac{[\ce{NO2}]^2}{[\ce{N2O4}]} = \dfrac{(2x)^2}{1-x} = \dfrac{4x^2}{1-x} = 2{,}0$.
  \item $4x^2 = 2(1-x) \Rightarrow 4x^2 + 2x - 2 = 0 \Rightarrow 2x^2 + x - 1 = 0$. Discriminant
        $\Delta = 1 + 8 = 9$, $\sqrt{\Delta}=3$ :
        \[ x = \frac{-1 \pm 3}{4} \;\Rightarrow\; x = \rep{0{,}5} \quad(\text{la racine } x=-1 \text{ est rejetée}). \]
  \item $[\ce{N2O4}] = 1-0{,}5 = \rep{\SI{0.5}{\mol\per\litre}}$ ;\;
        $[\ce{NO2}] = 2\times0{,}5 = \rep{\SI{1.0}{\mol\per\litre}}$.
  \item Vérification : $K = \dfrac{(1{,}0)^2}{0{,}5} = 2{,}0$. \checkmark
\end{enumerate}
\end{correction}

\begin{correction}[équilibre symétrique : l'astuce de la racine carrée]
\begin{enumerate}[label=\alph*),leftmargin=2em,topsep=2pt,itemsep=2pt]
  \item Tableau (\si{\mol\per\litre}) :
  \begin{center}\small
  \begin{tabular}{@{}lccccc@{}}
  \toprule
   & $\ce{H2}$ & $+$ & $\ce{I2}$ & $\rightleftharpoons$ & $2\,\ce{HI}$ \\
  \midrule
  initial   & $2{,}0$ & & $2{,}0$ & & $0$ \\
  équilibre & $2-x$   & & $2-x$   & & $2x$ \\
  \bottomrule
  \end{tabular}
  \end{center}
  \item $K = \dfrac{(2x)^2}{(2-x)(2-x)} = \left(\dfrac{2x}{2-x}\right)^2 = 16$. Les deux membres étant des
        \emph{carrés de quantités positives} ($x>0$ et $2-x>0$), on peut prendre la racine carrée.
  \item $\dfrac{2x}{2-x} = \sqrt{16} = 4 \Rightarrow 2x = 4(2-x) = 8 - 4x \Rightarrow 6x = 8 \Rightarrow
        x = \textcolor{primarydark}{\boldsymbol{\tfrac{4}{3}}} \approx 1{,}33$.
  \item $[\ce{HI}] = 2x = \dfrac{8}{3} \approx \rep{\SI{2.67}{\mol\per\litre}}$.
        (On vérifie $[\ce{H2}]=[\ce{I2}]=2-\tfrac43=\tfrac23\approx\SI{0.67}{\mol\per\litre}$.)
\end{enumerate}
\end{correction}

\begin{correction}[du taux de décomposition à la constante $K$]
\begin{enumerate}[label=\alph*),leftmargin=2em,topsep=2pt,itemsep=2pt]
  \item $n_0 = \dfrac{12{,}75}{51} = \rep{\SI{0.25}{\mol}}$.
  \item Décomposé : $0{,}25 \times 0{,}20 = \SI{0.05}{\mol}$. Chaque mole de $\ce{NH4HS}$ donne $1$ mole de
        $\ce{NH3}$ et $1$ mole de $\ce{H2S}$ : $n(\ce{NH3}) = n(\ce{H2S}) = \rep{\SI{0.05}{\mol}}$.
  \item $[\ce{NH3}] = [\ce{H2S}] = \dfrac{0{,}05}{2{,}00} = \rep{\SI{0.025}{\mol\per\litre}}$.
  \item Le solide $\ce{NH4HS}$ n'apparaît pas : $K_c = [\ce{NH3}]\,[\ce{H2S}] = 0{,}025 \times 0{,}025 =
        \rep{\num{6.25e-4}}$.
\end{enumerate}
\end{correction}

\begin{correction}[optimiser un équilibre : tous les leviers réunis]
\begin{enumerate}[label=\alph*),leftmargin=2em,topsep=2pt,itemsep=2pt]
  \item Le carbone solide est exclu : $K = \dfrac{[\ce{CO}]^2}{[\ce{CO2}]}$.
  \item Gauche : $1$ mol de gaz ($\ce{CO2}$) ; droite : $2$ mol ($2\,\ce{CO}$). Augmenter le volume (baisser
        la pression) favorise le côté à \emph{plus} de gaz, soit la \rep{droite} $\to$ plus de $\ce{CO}$. \oui
  \item Réaction endothermique : augmenter $T$ déplace vers les produits (\rep{droite}) et \rep{augmente} $K$
        $\to$ plus de $\ce{CO}$. \oui
  \item (i) Retirer du $\ce{CO2}$ (réactif) : déplace vers la gauche $\to$ \emph{moins} de $\ce{CO}$. \non \quad
        (ii) Catalyseur : aucun déplacement (agit seulement sur la vitesse). \non
  \item Pour augmenter le rendement en $\ce{CO}$ : \rep{augmenter la température} et \rep{augmenter le volume}
        (baisser la pression).
\end{enumerate}
\end{correction}

\end{document}
